\section{結言}
この日の話題を横断すると、生成AIを実際に使う段階での性能評価、アクセス条件、運用費用に関心が移っている。モデルの速度だけでなくタスクの完了品質を測るRealmarkの提案と、エージェントのループによる費用リスクの指摘は、実用性を測る観点でつながっている。

ただし、観測された話題の多くは利用者の感想や二次情報であり、NVIDIAの経営、Claudeの価格比較、Cursorのセキュリティ更新などには未確認事項が残る。Grok自身も検索結果が直近の投稿に偏り、観測窓全体の網羅性が限定的だったと評価している。本号は当日のX上の議論を記録したものであり、未確認の主張を確定情報とみなさないことが重要である。
